\documentclass[leqno]{article}
% Shared QBT/Soln pack preamble for resources.danieljsmith.org.
\usepackage{graphicx}
\usepackage{dsfont}
\usepackage{mathtools}
\usepackage{amssymb}
\usepackage{amsmath}
\usepackage{amsthm}
\usepackage{float}
\usepackage{ragged2e}
\usepackage{tensor}
\usepackage{fancyhdr}
\usepackage{empheq}
\usepackage[most]{tcolorbox}
\usepackage{enumitem}
\usepackage{dirtytalk}
\usepackage{marginnote}
\usepackage{microtype}
\usepackage[
  a4paper,
  left=0.8in,
  right=1in,
  top=1in,
  bottom=0.5in,
  headheight=14pt,
  headsep=0.4in
]{geometry}
\setlength{\parindent}{0cm}

% TikZ / PGFPlots
\usepackage{pgfplots}
\pgfplotsset{compat=1.18}
\usepgfplotslibrary{fillbetween, polar}
\usetikzlibrary{calc, positioning}
\usetikzlibrary{angles, quotes}
\usetikzlibrary{arrows.meta}
\usetikzlibrary{decorations.pathreplacing, decorations.markings}
\usetikzlibrary{patterns}
\usetikzlibrary{intersections}

% Links and contents
\usepackage[colorlinks=true,linkcolor=black,citecolor=magenta,urlcolor=black]{hyperref}%
\urlstyle{same}

% Pack macros
\RenewDocumentCommand{\marks}{o m}{%
  \IfValueTF{#1}%
    {\marginnote{\raggedright\textbf{[#2]}}[#1]}%
    {\marginnote{\raggedright\textbf{[#2]}}}%
}

\newcommand{\vspacerule}[2]{%
  \par\vspace*{#1}%
  \noindent\hrulefill\par
  \vspace*{#2}%
}

\definecolor{scarlet}{RGB}{205, 40, 75}

% Worked-solution box, adapted from the TMUA worked-solutions box. A white breakable box with a left scarlet rule and no surrounding
% frame or tint. A part that breaks fills the rest of its page, so the rule
% runs to the bottom margin; the unbroken or last part ends in a short
% horizontal foot rule so a solution visibly stops. Unlike the TMUA box there
% is no answer label and no extra \parskip: these solutions mix \\ line breaks
% with blank lines, so paragraph skips would make the spacing uneven. The box
% does not grow to the left, so the rule sits at the text edge; growing to the
% right keeps the usable line width close to the old tcolorbox Solution box.
\newtcolorbox{djsSolution}{%
  enhanced, breakable, frame hidden, sharp corners, colback=white,
  height fixed for=first and middle, height fill,
  extras unbroken and last={natural height},
  borderline west={2.2pt}{0pt}{scarlet},
  left=11pt, right=0pt, top=5pt, bottom=13pt,
  before skip=13pt, after skip=4pt,
  grow to right by=10mm,
  overlay unbroken and last={%
    \fill[scarlet] (frame.south west)
      rectangle ([xshift=3.5cm, yshift=2.2pt]frame.south west);},
  before upper={{\color{scarlet}\bfseries Solution}\par\vspace{4pt}}}

% Optional smaller figure inside a solution box. \scalebox shrinks node text
% with the drawing. Not applied to existing figures.
\newcommand{\SolnFigureScale}{0.85}
\newsavebox{\djsSolnFigureBox}
\newenvironment{djsSolutionFigure}
  {\par\vspace{4pt}\begin{center}\begin{lrbox}{\djsSolnFigureBox}}
  {\end{lrbox}\scalebox{\SolnFigureScale}{\usebox{\djsSolnFigureBox}}%
   \end{center}\vspace{2pt}\par}

\newcommand{\cxl}[2]{%
  \tikz[baseline=(X.base)]{%
    \node[inner sep=0pt, outer sep=0pt, text=#1] (X) {$#2$};
    \draw[#1, line width=0.4pt] (X.south west)--(X.north east);
  }%
}

% Page-style hook run at the contents page and each question's opening page.
% A no-op for QBT packs; \djsSolnHeader points it at the fancy header.
\newcommand{\djsQuestionPageStyle}{}

\newcommand{\questionitem}{%
  \djsQuestionPageStyle
  \item\phantomsection
  \addcontentsline{toc}{section}{Question \number\value{enumi}}%
}

\renewcommand{\contentsname}{Questions}

\newcommand{\djsPackHeader}[1]{%
  \pagestyle{fancy}%
  \fancyhead{}%
  \fancyfoot{}%
  \renewcommand{\headrulewidth}{0.9pt}%
  \lhead{#1}%
  \chead{}%
  \rhead{\href{https://danieljsmith.org}{danieljsmith.org}}%
}

\newcommand{\djsQbtHeader}[1]{%
  \djsPackHeader{#1}%
}

% Solutions packs show the header only on the contents page and each
% question's opening page, so it never appears part-way through a solution.
% \thispagestyle marks the next page shipped out, which is the question's
% page because every solution question starts after a \newpage.
\newcommand{\djsSolnHeader}[1]{%
  \djsPackHeader{\textit{Solutions} - #1}%
  \pagestyle{empty}%
  \renewcommand{\djsQuestionPageStyle}{\thispagestyle{fancy}}%
}

\newcommand{\djsFrontMatter}{%
  \djsQuestionPageStyle
  \vspace*{20pt}%
  \tableofcontents
  \newpage
}


\djsQbtHeader{Dimensional Analysis}

\begin{document}
\djsFrontMatter
\begin{enumerate}[
  label=\textbf{\arabic*.},
  leftmargin=*,
  topsep=0pt,
  partopsep=0pt,
  parsep=0pt
]

% ---- Question 1 ----
\questionitem
% [Source: _QBT__Dimensional_Analysis, Q1]

\begin{enumerate}[label=\textbf{(\alph*)}]
  \item State the dimensions of force. \marks{1} \\
  \item The braking distance of a vehicle can be modelled by the formula
  \[
  d = m^a u^b R^c
  \] 
  where, \\

  $m$ is the mass of the vehicle in $\mathrm{kg}$

  $u$ is the initial speed in $\mathrm{m\,s^{-1}}$

  $R$ is the constant braking force in $\mathrm{N}$ \\
  
  Use dimensional analysis to find the values of $a$, $b$ and $c$. \marks{4} \\
\end{enumerate}

\newpage

% ---- Question 2 ----
\questionitem
% [Source: _QBT__Dimensional_Analysis, Q2]

A small satellite moves in a circular orbit of radius $r$ about the centre of a much more massive planet of mass $M$. It is suggested that the period $T$ of the orbit depends only on $G$, $M$ and $r$, and that \\
\[
T = kG^{\alpha} M^{\beta} r^{\gamma}
\] \\
where $k$ is a dimensionless constant. \\

It is given that
\[
G = 6.67 \times 10^{-11}\,\mathrm{m}^3\,\mathrm{kg}^{-1}\,\mathrm{s}^{-2}
\] \\
\begin{enumerate}[label=\textbf{(\alph*)}, resume]
  \item Write down the dimensions of $G$. \marks{1} \\
  \item By using dimensional analysis, determine the values of $\alpha$, $\beta$ and $\gamma$. \marks{3} \\
\end{enumerate}

\newpage

% ---- Question 3 ----
\questionitem
% [Source: _QBT__Dimensional_Analysis, Q3]

\begin{figure}[H]
    \centering
\begin{tikzpicture}[scale=1.1]
\def\Height{3.0}
\def\Range{8.2}
\def\Angle{38}
\def\VelLen{2.3}
\def\DimDrop{0.8}
\def\DimOff{0.55}
\pgfmathsetmacro{\TanA}{tan(\Angle)}
\pgfmathsetmacro{\Coeff}{(\Height+\TanA*\Range)/(\Range*\Range)}

\coordinate (B) at (0,0);
\coordinate (O) at (0,\Height);
\coordinate (A) at (\Range,0);
\coordinate (Href) at ($(O)+(1.6,0)$);
\coordinate (V) at ($(O)+(\Angle:\VelLen)$);
\coordinate (R1) at ($(B)+(0,-\DimDrop)$);
\coordinate (R2) at ($(A)+(0,-\DimDrop)$);
\coordinate (h1) at ($(B)+(-\DimOff,0)$);
\coordinate (h2) at ($(O)+(-\DimOff,0)$);
\coordinate (h3) at ($(O)+(2*\DimOff,0)$);

\draw[thick] (-1.2,0) -- (\Range+1.2,0);
\draw[thick] (B) -- (O);
\draw[thick, black, domain=0:\Range, samples=200]
  plot (\x, {\Height+\TanA*\x-\Coeff*\x*\x});

\draw[-{Stealth[length=3mm]}, thick] (O) -- (V)
  node[pos=0.65, above] {$v$};

\pic["$\theta$", draw, angle radius=7mm, angle eccentricity=1.35]
  {angle=Href--O--V};

\draw[dashed] (B) -- (h1);
\draw[dashed] (O) -- (h3);
\draw[{Stealth[length=3mm]}-{Stealth[length=3mm]}, thick] (h1) -- (h2)
  node[midway, left] {$h$};

\draw[{Stealth[length=3mm]}-{Stealth[length=3mm]}, thick] (R1) -- (R2)
  node[midway, below] {$R$};

\fill[black] (A) circle (1.2pt);

\node[above left] at (O) {$O$};
\node[below left] at (B) {$B$};
\node[below right] at (A) {$A$};
\end{tikzpicture}
\end{figure}


A particle is projected with speed $v$ from a point $O$ at the top of a cliff of height $h$ above horizontal ground. The angle of projection is $\theta$ above the horizontal. The particle lands at the point $A$ on the ground, and $B$ is the foot of the cliff, so that $BA = R$, as shown in the diagram. \\

You are given that
\[
R = \frac{v^\alpha \cos\theta}{g^\beta}\left(v^\gamma\sin\theta + \sqrt{v^\delta\sin^2\theta + 2gh}\right)
\] \\
Use dimensional analysis to find $\alpha$, $\beta$, $\gamma$ and $\delta$. \marks{4} \\

\newpage

% ---- Question 4 ----
\questionitem
% [Source: _QBT__Dimensional_Analysis, Q4]

The increase in pressure $p$ pascals between the surface and the bottom of a liquid of depth $h$ metres and density $\rho\,\mathrm{kg\,m^{-3}}$ is modelled by
\[
p = k\rho^{\alpha}h^{\beta}g
\] \\
where $g\,\mathrm{m\,s^{-2}}$ is the acceleration due to gravity and $k$ is a dimensionless constant. \\
\begin{enumerate}[label=\textbf{(\alph*)}]
  \item State the dimensions of $p$. \marks{1} \\
  \item Use dimensional analysis to find the value of $\alpha$ and the value of $\beta$. \marks{3} \\
\end{enumerate}

\newpage

% ---- Question 5 ----
\questionitem
% [Source: _QBT__Dimensional_Analysis, Q5]

The force, $F$, exerted on a flat plate by a water jet is thought to depend on the density of the water, $\rho$, the cross-sectional area of the jet, $A$, and the speed of the water, $v$. You are given that $\rho$ is measured in $\mathrm{kg\,m^{-3}}$. \\
\begin{enumerate}[label=\textbf{(\alph*)}]
  \item Find the dimensions of density $\rho$. \marks{2} \\
\end{enumerate}
A model of the force suggests the relationship
\[
F = \frac12 \rho^\alpha A^\beta v^\gamma
\] \\
\begin{enumerate}[label=\textbf{(\alph*)}, resume]
  \item Explain whether or not it is possible to use dimensional analysis to verify that the constant $\frac12$ is correct. \marks{1} \\
  \item Use dimensional analysis to find the values of $\alpha$, $\beta$ and $\gamma$. \marks{4} \\
\end{enumerate}

\newpage

% ---- Question 6 ----
\questionitem
% [Source: _QBT__Dimensional_Analysis, Q6]

A motor pulls a crate of mass $m$ along a rough horizontal floor. The equation below is suggested to model the pulling force exerted by the motor, $F$, while the crate moves in a straight line. \\
\[
F = k_1 m v\dfrac{\mathrm{d}v}{\mathrm{d}x} + k_2 \mu mg + k_3 \dfrac{Q}{x}
\] \\
in which, \\

$\bullet$ $v$ is the speed of the crate when it has travelled distance $x$

$\bullet$ $\mu$ is the coefficient of friction between the crate and the floor

$\bullet$ $g$ is the acceleration due to gravity

$\bullet$ $Q$ is the total thermal energy transferred to the surroundings after the crate has moved distance $x$

$\bullet$ $k_1$, $k_2$ and $k_3$ are dimensionless constants \\

Determine whether the equation is dimensionally consistent. \marks{6} \\

\newpage


% ---- Question 7 ----
\questionitem
% [Source: _QBT__Dimensional_Analysis, Q7]

The speed of a particle falling vertically through a resistive medium is modelled by the formula
\[
v=\dfrac{mg}{k}\left(1-e^{-kt/m}\right)
\] 
where, \\

$v$ is the speed of the particle

$m$ is the mass of the particle

$g$ is the acceleration due to gravity

$t$ is the time since the particle was released

$k$ is a constant \\

Use dimensional analysis to determine the dimensions of $k$. \marks{4} \\

\newpage

% ---- Question 8 ----
\questionitem
% [Source: _QBT__Dimensional_Analysis, Q8]

A comet of mass $m$ measured in $\mathrm{kg}$ follows an elliptical orbit around a planet. The semi-major axis of the orbit is $a$ measured in metres and the orbital period is $T$ measured in seconds. The eccentricity $e$ of the orbit is dimensionless. \\

A constant $\mu$ associated with the planet is defined by
\[
T^2 = \frac{4\pi^2 a^3}{\mu}
\] \\
and the angular momentum of the comet is given by
\[
H = m\sqrt{\mu a(1-e^2)}
\] \\
By using dimensional analysis, find the dimensions of:

\begin{enumerate}[label=\textbf{(\alph*)}]
  \item $\mu$ \marks{3}
  \item $H$ \marks{3}
\end{enumerate}

\newpage

% ---- Question 9 ----
\questionitem
% [Source: _QBT__Dimensional_Analysis, Q9]

A small ball of mass $m\,\mathrm{kg}$ is projected vertically upwards with speed $u\,\mathrm{m\,s^{-1}}$. During the motion, when its speed is $v\,\mathrm{m\,s^{-1}}$, it experiences an air resistance of magnitude $kv^2\,\mathrm{N}$, where $k$ is a constant. \\
\begin{enumerate}[label=\textbf{(\alph*)}]
  \item Determine the dimensions of $k$. \marks{3} \\
\end{enumerate}
The greatest height reached by the ball is $H\,\mathrm{m}$. A formula approximating $H$ is
\[
H = \dfrac{u^2}{2g} - \dfrac{ku^4}{4mg^2} + \dfrac{1}{6}k^2 m^\alpha u^\beta g^\gamma
\] \\
where $g\,\mathrm{m\,s^{-2}}$ is the acceleration due to gravity. \\
\begin{enumerate}[label=\textbf{(\alph*)}, resume]
  \item Use dimensional analysis to find $\alpha$, $\beta$ and $\gamma$. \marks{4} \\
\end{enumerate}

\newpage 

% ---- Question 10 ----
\questionitem
% [Source: _QBT__Dimensional_Analysis, Q10]

A small probe is launched vertically upwards from the surface of a spherical moon. The radius of the moon is $R$ metres, and the probe is at a distance $r$ metres from the centre of the moon. The gravitational field strength $g$ at distance $r$ is modelled by
\[
g = \frac{\Gamma}{r^2}
\]
where $\Gamma$ is a constant associated with the moon. \\

The probe has initial speed $u$ at the surface and speed $v$ when it reaches distance $r$. These speeds are related by
\[
v^2 = u^2 - 2\Gamma\left(\frac{1}{R}-\frac{1}{r}\right)
\]
The work done against gravity in moving a probe of mass $m$ from the surface to distance $r$ is
\[
W = m\Gamma\left(\frac{1}{R}-\frac{1}{r}\right)
\]

\begin{enumerate}[label=\textbf{(\alph*)}]
  \item Find the dimensions of $\Gamma$. \marks{2}
  \item Use dimensional analysis to show that the equation for $v^2$ is dimensionally consistent. \marks{2}
  \item Find the dimensions and SI units of $W$. \marks{2}
\end{enumerate}


\end{enumerate}
\end{document}
